\subsection{完全格序空间}

定义 2. --- 一个 Riesz 空间 E 称为完全格序的，如果 E 的每个有上界的非空子集在 E 中有一个上确界。

直接可知，在完全格序空间 E 中，每个有下界的非空子集在 E 中有一个下确界。

例. --- 1) 若 A 是任意集，则定义在 A 上的实值函数所成的空间 \(R^{A}\) 是完全格序的，因为有上界的族在 \(R^{A}\) 中的上确界就是它的上包络 (GT, IV, §5, No. 5)。

\begin{enumerate}
\def\labelenumi{\arabic{enumi})}
\setcounter{enumi}{1}
\tightlist
\item
  设 F 是任意集；F 上有界实值函数所成的空间 \(\mathcal{B}(\mathrm{F})\) ，赋以由 \(R^{F}\) 的序结构诱导的序结构，是完全格序的。然而，若 F 是拓扑空间，则 F 上连续实值函数所成的空间 \(\mathcal{C}(\mathrm{F})\) （赋以由 \(R^{F}\) 的序结构诱导的序结构）是一个 Riesz 空间，但一般说来不是完全格序的（参见习题 13）。例如考虑 \(F = R\) 的情形；设 I 是区间 {]}0,1{[}，\(\varphi_{I}\) 是 I 的特征函数，又设 H 是连续函数 \(x(t)\) （满足 \(x \leqslant \varphi_{I}\) ）所成的集；显然 H 在 \(\mathcal{C}(\mathrm{F})\) 中有上界。函数 \(\varphi_{I}\) 是诸 \(x \in H\) 的上包络，但它不是它们在 \(\mathcal{C}(\mathrm{F})\) 中的上确界，因为 \(\varphi_{I}\) 下半连续而不连续。我们来证明，事实上 H 在 \(\mathcal{C}(\mathrm{F})\) 中没有上确界；只需证明：若 u 是使 \(u \geqslant \varphi_{I}\) 的连续函数，则存在连续函数 \(v \neq u\) 使 \(u \geqslant v \geqslant \varphi_{I}\) 。而 \(u(0) \geqslant 1\) ，因此存在数 \(\alpha > 0\) 使 \(u(t) > 0\) 对 \(-\alpha \leqslant t \leqslant 0\) 成立；若 w 是一个在区间 {]}- \(\alpha,0[\) 之外为零且在该区间上满足 \(0 < w(t) < u(t)\) 的连续函数，则函数 v = u - w 即满足要求。
\end{enumerate}

命题 1. --- 为使一个序向量空间 E 是完全格序的，必要且充分的是：E 是一个 Riesz 空间，且满足以下两个条件之一：

\begin{enumerate}
\def\labelenumi{\alph{enumi})}
\item
  每个非空子集 A（由 E 中元素 \(\geqslant0\) 组成），若有上界且依关系 \(\leqslant\) 有向，则在 E 中有一个上确界；
\item
  每个非空子集 A（由 E 中元素 \(\geqslant 0\) 组成且依关系 \(\geqslant\) 有向），在 E 中有一个下确界。
\end{enumerate}

这些条件显然是必要的。反过来，设 E 是满足条件 a) 的 Riesz 空间。设 B 是 E 的一个有上界的非空子集；由 B 的有限子集的上确界组成的集 C 依关系 \(\leqslant\) 有向；设 a 是它的一个元素，\(C_{a}\) 是由 \(x \in C\) 中那些 \(\geqslant a\) 者所成的集；若我们能证明 \(C_{a}\) 有上确界，则它也是 B 的上确界。而 \(C_{a} - a\) 是一个由元素 \(\geqslant 0\) 组成的集，有上界且依关系 \(\leqslant\) 有向；因此它有上确界 b，从而 \(a + b\) 是 \(C_{a}\) 的上确界。

另一方面，条件 \(b)\) 蕴涵 \(a)\) ：因为若 \(F\) 是一个非空集，由元素 \(\geqslant 0\) （取自 \(E\) ）组成，有上界且依 \(\leqslant\) 有向，又若 \(c\) 是 \(F\) 的一个上界，则 \(c - F\) 是一个由元素 \(\geqslant 0\) 组成且依 \(\geqslant\) 有向的集；若它有下确界 \(m\) ，则 \(c - m\) 是 \(F\) 的上确界。

命题 2. --- 设 E 是一个 Riesz 空间，赋以一个与其序向量空间结构相容的 Hausdorff 拓扑 (TVS, II, §2, No. 7)。若对 E 的每个有上界且依关系 ≤ 有向的集 H ⊂ E，H 的截断滤子都收敛，则 E 是完全格序的。

事实上，已知 H 的截断滤子的极限就是 H 在 E 中的上确界 (TVS, II, §2, No. 7, Prop. 18)。

